\documentclass[10pt,a4paper]{amsart}
\usepackage[margin=19mm]{geometry}
\usepackage{amsmath,amssymb,mathtools}
\usepackage[scr=rsfs,cal=euler]{mathalfa}
\usepackage{xfrac}
\usepackage[unicode,hidelinks,bookmarks=false]{hyperref}
\newcommand{\ie}{i.e.\ }
\newcommand{\cf}{cf.\ }
\newcommand{\resp}{respectively\ }
\newcommand{\Subsection}{\subsection}
\providecommand{\nobreakdash}{\nobreak\hbox{-}}
\setlength{\emergencystretch}{2em}
\multlinegap0pt
\usepackage{bm}
\usepackage{bbm}
\usepackage{dsfont}
\usepackage{xspace}
\usepackage{cancel}
\usepackage{mathtools}
\usepackage{amsthm}
\makeatletter
\@ifundefined{theorem}{\newtheorem{theorem}{Theorem}}{}
\@ifundefined{remark}{\newtheorem{remark}[theorem]{Remark}}{}
\makeatother
\def\DDD{\mathscr{D}}
\def\EE{{\mathcal E}}
\def\MM{{\mathcal M}}
\def\P{\vec{P}}
\def\Q{\vec{Q}}
\def\R{{\mathbb R}}
\def\TT{{\mathcal T}}
\def\Z{{\mathbb Z}}
\def\adown{a^{\downarrow}}
\def\al{\alpha}
\def\aup{a^{\uparrow}}
\def\dd{\delta}
\def\dist{{\mathscr{D}}}
\def\ee{{e}}
\def\ep{\varepsilon}
\def\f{\vec{f}}
\def\ga{\vec{\gamma}}
\def\g{\vec{g}}
\def\l1{\Lambda^{1}}
\def\lzero{\Lambda^0}
\def\xxi{\vec{\xi}}
\def\x{\vec{x}}
\newcommand\assump[1]{{\rm\textbf{#1}}}
\newcommand\bs[1]{{\boldsymbol{#1}}}
\title[A new construction for Melnikov chaos in piecewise-smooth pl]{A new construction for Melnikov chaos in piecewise-smooth planar systems}
\author{Alessandro Calamai, Matteo Franca, Michal Pospisil}
\date{Source: arXiv:2507.15543v1 (CC BY 4.0)}
\begin{document}
\maketitle

\noindent\emph{Source and licence.} A new construction for Melnikov chaos in piecewise-smooth planar systems. Version arXiv:2507.15543v1 (CC BY 4.0), distributed under the Creative Commons Attribution 4.0 International licence, as recorded in the arXiv licence metadata for this exact version and shown on the abstract page https://arxiv.org/abs/2507.15543v1. Full rights notice: licenses/arxiv-2507.15543-CC-BY-4.0.txt.\par\smallskip

\noindent\textit{Editorial scope.} Counted unit: the Theorem whose source label is \texttt{theoremkeybis} in the article (source0.tex lines 2336--2344, source label \texttt{theoremkeybis}), delivered together with the article's own complete original proof of it (source0.tex lines 2345--2383, one \texttt{proof} environment of 39 lines). No mathematics is written, repaired or re-derived: statement and proof are the article's own text line for line. Required context retained uncounted: minimal (lines 799--806); defbetaj (lines 828--834); isolated (lines 836--838); TandKnu (lines 852--857); TandKnunew (lines 871--876); main.weak1 (lines 938--974); veryweak (lines 1082--1094); theoremkey (lines 2203--2209); rMelni (lines 2248--2256). Every definition, display and label of those lines is retained unchanged. Omitted from this package and not redistributed: 1--798, 807--827, 835--835, 839--851, 858--870, 877--937, 975--1081, 1095--2202, 2210--2247, 2257--2335, 2384--3303. Nothing inside a retained excerpt is omitted or altered: every retained body is delivered exactly as the article's lines are written, so no omission receipt of any kind accompanies this package. Citations: the retained excerpts carry no citation command at all, so no bibliography entry of the article is transcribed and no reference is redistributed. Reference adaptation: none was needed and none was made; every reference inside the delivered unit resolves inside it, so no unresolved reference or citation is produced. Printed slips kept literal: no printed word, symbol, hypothesis or number of the counted statement or of its proof was altered. Layout and class: the article's own class is \texttt{article} with packages including amsthm, amsfonts, amssymb, amsmath, amscd, latexsym, mathrsfs, epsfig, inputenc, xy, cite, subcaption, xcolor; the wrapper is the lane's pinned \texttt{amsart} editorial wrapper at 10pt on A4, which restates the theorem-like environments the retained text needs and adds the article's own macro definitions that the retained text needs (\texttt{\textbackslash DDD}, \texttt{\textbackslash EE}, \texttt{\textbackslash MM}, \texttt{\textbackslash P}, \texttt{\textbackslash Q}, \texttt{\textbackslash R}, \texttt{\textbackslash TT}, \texttt{\textbackslash Z}, \texttt{\textbackslash adown}, \texttt{\textbackslash al}, \texttt{\textbackslash assump}, \texttt{\textbackslash aup}, \texttt{\textbackslash bs}, \texttt{\textbackslash dd}, \texttt{\textbackslash dist}, \texttt{\textbackslash ee}, \texttt{\textbackslash ep}, \texttt{\textbackslash f}, \texttt{\textbackslash g}, \texttt{\textbackslash ga}, \texttt{\textbackslash l}, \texttt{\textbackslash lzero}, \texttt{\textbackslash x}, \texttt{\textbackslash xxi}), with extra wrapper packages (bm, bbm, dsfont, xspace, cancel, mathtools). The delivered numbering is the unit's own. Assets: no retained span contains a figure environment, a graphics-inclusion command, a PDF-page inclusion, a table or a TikZ picture; no image file is copied into this package, the package declares no asset dependency and no image file is redistributed with it. Licence: version arXiv:2507.15543v1 of this article is distributed under the Creative Commons Attribution 4.0 International licence, as recorded in the arXiv licence metadata for this exact version and shown on the abstract page https://arxiv.org/abs/2507.15543v1. A full rights notice is delivered at licenses/arxiv-2507.15543-CC-BY-4.0.txt. Characters: every retained excerpt is pure ASCII, so the delivered engine, XeLaTeX with its default Latin Modern text font, reports no missing character.
\begin{equation}\label{minimal}
\begin{split}
 & |\MM(\aup_j)|=|\MM(\adown_j)|= \dd \bar{c}, \qquad \MM(\aup_j) \MM(\adown_j) <0, \\ & \MM(\tau) \ne 0 \quad \forall \tau \in [\aup_j,
  \adown_j] \setminus [T_{2j}-\lzero, T_{2j}+\lzero],\\
  & \textrm{$0<\adown_j-\aup_j \le \l1$,
for any $j \in \Z$.}
  \end{split}
\end{equation}

  \begin{equation}\label{defbetaj}
\begin{split}
    &  |\MM(\aup_j)|=|\MM(\adown_j)|= \dd \bar{c}, \qquad \MM(\aup_j) \MM(\adown_j) <0, \\
     & \MM(\tau) \ne 0 \qquad \forall \tau \in [\aup_j,
  \adown_j] \setminus \{T_{2j} \}.
\end{split}
\end{equation}

\begin{equation}\label{isolated}
  |\MM(T_{2j}+ h)| \ge \omega_M(|h|)  \quad  \textrm{for any }  (T_{2j}+ h)\in [\aup_j, \adown_j]\, , \quad h \ne 0,
\end{equation}

\begin{equation}\label{TandKnu}
\begin{split}
  &  \MM(T_{2j})=0, \quad T_0 \in [b_0,b_1],\\
& T_{j+1}-T_j > \l1 + K_0 (1+\nu) |\ln(\ep)| ,   \quad \textrm{ for any $j \in \Z$}
  \end{split}
\end{equation}

\begin{equation}\label{TandKnunew}
\begin{split}
  &  \MM(T_{2j})=0, \quad T_0 \in [b_0,b_1],\\
& T_{j+1}-T_j >    \max\{ B_{j+1}; B_{j}\}+ K_0 (1+\nu) |\ln(\ep)| ,   \textrm{ for any $j \in \Z$.}
  \end{split}
\end{equation}

\begin{theorem}\label{main.weak1}
	Assume that   $\f^{\pm}$ and $\g$ are $C^r$ with  $r > 1$,   and that \assump{F0}, \assump{F1}, \assump{F2}, \assump{K} and \assump{G} hold true;
assume further    \assump{P1}.
Then we can choose $\ep_0$ small enough so that for any $0< \ep \le \ep_0$
and any increasing sequence $\mathcal{T}=(T_j)$, $j \in \Z$,  satisfying  \eqref{TandKnu} and \eqref{minimal},
 there are $c^*>0$ (independent of $\ep>0$) and  a continuous increasing function $\omega$ satisfying
  $\omega(0)=0$
so that we have the following.
For any $\ee=(\ee_j) \in \EE$,
there is a compact set $\bar{X}(\ee, \mathcal{T})$ and a sequence $(\al_j^{\ee}(\ep))$, $j \in \Z$,
such that $|\al_j^{\ee}(\ep)| \le \omega(\ep)$ and for any $\xxi(\ee) \in \bar{X}(\ee, \mathcal{T})$
  the trajectory $\x(t,T_0 + \al_0^{\ee}(\ep); \xxi(\ee))$ has the property $\mathbf{C_e}$, i.e.
  \begin{description}
  \item[$\mathbf{C_e}$]
	if $\ee_j=1$, then
	\begin{equation}\label{ej=1w}
	\begin{split}
	&   \|\x(t,T_0+ \al_0^{\ee}(\ep); \xxi(\ee) )-\ga(t-T_{2j}-\al_j^{\ee}(\ep)) \| \le c^*\ep  \\
 & \textrm{when $t \in [T_{2j-1}, T_{2j+1} ]$ }
	\end{split}
	\end{equation}
	while if $\ee_j=0$, we have
	 \begin{equation}\label{ej=0w}
\begin{split}
	&	\|\x(t,T_0+ \al_0^{\ee}(\ep); \xxi(\ee))  \| \le c^*\ep \\
 & \textrm{when $t \in [T_{2j-1}, T_{2j+1} ]$ }
	\end{split}
	\end{equation}
	for any $j \in \Z$.
\end{description}
Moreover, we have the following estimates for $\alpha_j^{\ee}(\ep)$:
\begin{itemize}
  \item  Assume \eqref{minimal}, then $|\alpha_j^{\ee}(\ep)| \le  \l1$ for any $j \in \Z$.
  \item Assume \eqref{isolated}, then there is   a continuous increasing function $\omega_{\alpha}(\ep)$ (independent of $\ee$) satisfying
  $\omega_{\alpha}(0)=0$, such that $|\alpha_j^{\ee}(\ep)| \le \omega_{\alpha}(\ep)$ for any $j \in \Z$.
\end{itemize}
\end{theorem}

\begin{theorem}\label{veryweak}
	Assume that   $\f^{\pm}$ and $\g$ are $C^r$ with  $r > 1$,   and that \assump{F0}, \assump{F1}, \assump{F2}, \assump{K} and \assump{G} hold true;
assume further    \assump{P1}.
Then we can choose $\ep_0$ small enough so that for any $0< \ep \le \ep_0$
and any increasing sequence $\mathcal{T}=(T_j)$, $j \in \Z$,  satisfying  \eqref{TandKnunew},
 there is $c^*>0$ (independent of $\ep>0$ and  of the sequence $(T_j)$)
so that we have the following.
For any $\ee=(\ee_j) \in \EE$,
there is a compact set $\bar{X}(\ee, \mathcal{T})$ and a sequence $(\al_j^{\ee}(\ep))$, $j \in \Z$,
such that $|\al_j^{\ee}(\ep)| \le B_{2j}$ and for any $\xxi(\ee) \in \bar{X}(\ee, \mathcal{T})$
  the trajectory $\x(t,T_0+ \al_0^{\ee}(\ep); \xxi(\ee))$ has the property
  $\mathbf{C_e}$.
\end{theorem}

\begin{theorem}\label{theoremkey}
  Assume that  the hypotheses of Theorem  \ref{veryweak} hold. Let $\ee \in \EE$ be a sequence satisfying
   $\ee_0=1$. Then there is a compact  connected  set $X(\ee, \TT)\subset L^0$ and a sequence $(\al_j^{\ee}(\ep))$, $j \in
   \Z$,
   such that if $\xxi \in X(\ee, \TT)$,
   then $\x(t,T_0+\al_0^\ee(\ep); \xxi)$ has property  $\bs{C_{\ee}}$.  Further $|\alpha_0| \le B_0$.
\end{theorem}

\begin{remark}\label{rMelni}
	Assume that $\f^{\pm}$ and $\g$ are $C^r$, $r>1$.
Observe that if $0<\ep \ll 1$ then,  for any fixed $\tau \in \R$,
	there are $c>0$ (independent of $\ep$ and $\tau$) and a monotone increasing and continuous function
$\bar{\omega}(\ep)$ such that $\bar{\omega}(0)=0$ and
	$$\dist(\P_s(\tau),\P_u(\tau))= c \ep \MM(\tau)+ \ep  \omega(\tau,\ep) \, , \qquad  |\omega(\tau,\ep)| \le
\bar{\omega}(\ep).$$
Further, if $r \ge 2$, we can find $C_M>0$ such that $\bar{\omega}(\ep) \le C_M \ep$.
\end{remark}

\begin{theorem}\label{theoremkeybis}
  Assume that  the hypotheses of Theorem \ref{main.weak1} hold;   in particular let $\TT$ be a sequence satisfying \eqref{TandKnu}
  and \eqref{minimal}.
   Let $\ee \in \EE$ be a sequence satisfying
   $\ee_0=1$. Then there is a compact  connected  set $X(\ee, \TT)\subset L^0$  and a sequence $(\al_j^{\ee}(\ep))$, $j \in
   \Z$,
   such that if $\xxi \in X(\ee, \TT)$,
   then $\x(t,T_0+\al_0^\ee(\ep); \xxi)$ has property  $\bs{C_{\ee}}$.  Further $|\alpha_0| \le \l1$.
\end{theorem}

\begin{proof}
  We just need to adapt the argument of the proof of Theorem \ref{theoremkey}.





  We claim that
\\ $\bullet$ \emph{there is}
$(d_{\star}^{\Lambda},\tau_{\star}^{\Lambda}) \in (\hat{S}^{\Lambda,-}(\TT^-, \ee^-) \cap  \tilde{S}^{\Lambda,+}(\TT^+, \ee^+))$.
\\
In fact let
$d^{\pm}_{\uparrow} \in \tilde{J}_{\pm\infty}(\aup_0)$, $d^{\pm}_{\downarrow} \in \tilde{J}_{\pm\infty}(\adown_0)$,
 and  $\hat{d}^-_{\uparrow}= d^-_{\uparrow}+ \DDD(\P_u(\aup_0), \P_s(\aup_0)) \in \hat{J}_{-\infty}(\aup_0)$,
 $\hat{d}^-_{\downarrow}= d^-_{\downarrow}+ \DDD(\P_u(\adown_0), \P_s(\adown_0)) \in \hat{J}_{-\infty}(\adown_0)$;
from \eqref{minimal} and \eqref{defbetaj} we can assume  $\MM(\aup_0)=-\bar{c} \delta<0<\bar{c} \delta =\MM(\adown_0)$
(or we have opposite sign conditions but the argument goes through without significant changes);
hence using Remark \ref{rMelni}, for any $0<\ep \le \ep_0$ we find
\begin{equation}\label{d+id-ibis}
  \begin{split}
     d^{+}_{\uparrow}-\hat{d}^{-}_{\uparrow} \le  &   \frac{c}{2} \MM(\aup_0) \ep +\ep^{1+\bar{\nu} } \le    \frac{c}{4} \MM(\aup_0) \ep =  -
     \frac{c\bar{c}}{4} \delta \ep , \\
     d^{+}_{\downarrow}-\hat{d}^{-}_{\downarrow}\ge  &  \frac{c}{2} \MM(\adown_0)\ep  -\ep^{1+\bar{\nu} } \ge   \frac{c}{4} \MM(\adown_0) \ep
     =
     \frac{c\bar{c}}{4}  \delta\ep ,
  \end{split}
\end{equation}
 with $\bar{\nu}= \frac{1-\underline{\sigma}+\nu}{\underline{\sigma}}>1$.
Then, arguing as in the proof of Theorem \ref{theoremkey}, we prove the claim.

Now let $d^{\Lambda,-}_\star=d_\star^{\Lambda}- \DDD(\P_u(\tau_\star^{\Lambda}), \P_s(\tau_\star^{\Lambda}))$; then
$(d^{\Lambda,-}_\star,\tau_\star^{\Lambda}) \in
\tilde{S}^{\Lambda,-}=\tilde{S}^{\Lambda,-}(\TT^-, \ee^-)$ and
$$\xxi_\star = \Q_s(d_\star^{\Lambda},\tau_\star^{\Lambda})= \Q_u(d^{\Lambda,-}_\star,\tau_\star^{\Lambda})$$
is such that $\x(t,\tau_\star^{\Lambda}; \xxi_\star)$ has property $\bs{C_{\ee}}$ for any $t \in \R$ or, in other words,
$\x(t,T_0+\alpha_0,\xxi_\star)$ with $\alpha_0=\tau_\star^{\Lambda}-T_0$ has property $\bs{C_{\ee}}$ for any $t \in \R$.

Further by construction   $ |\alpha_0| = |\tau_\star^{\Lambda}-T_0| \le \adown_0-\aup_0 \leq \l1$, so the theorem is proved.
\end{proof}

\end{document}
